% ---
% id: receiver/def:port-dilation
% kind: def
% title: "Port dilation of a receiver instrument"
% summary: >-
%   A port dilation is an isometry from the declared inputs into a finite mode
%   set with disjoint ports that reproduces every port Gram. Any such map
%   carries the click statistics as the unitary dilation of the physical model
%   does, and it may have far fewer modes.
% prominence: standard
% status: implemented
% uses: [def:instrument, def:gram, def:fock]
% mentions: [def:physmodel, thm:port-grams, thm:factor-blocks]
% assumes: []
% refs: []
% provenance: newly-derived
% implements: certiqs_key.receiver.factor:PortDilation
% lean:
% notes: "Stated on the port Grams alone, so that a declaration of port Grams or of an instrument matrix and a declared element graph all have port dilations; the map is an isometry from the declared inputs and need not extend the instrument matrix on the other modes."
% ---

The \tref{measurement isometry} of Definition~\ref{def:physmodel} applies a unitary
on the modes and an environment to the declared inputs, and
Theorem~\ref{thm:port-grams}\,(iii) and~(v) show that the click statistics
read of it only the amplitudes it sends from the inputs to each port. The
definition names the least that a map from the inputs must satisfy to carry
the same statistics, so that a smaller isometry, into the port modes and as
many loss modes as the instrument's loss needs, can stand in for the unitary
dilation: Theorem~\ref{thm:factor-blocks} reads every photon-number block
through such a map. The restriction $U\iota$ of a unitary dilation $U$ to the
declared inputs, with the ports of the instrument, is one, since
$(U\iota)^{\dagger}U\iota=\id_{\Inp}$ and the rows of $U\iota$ on a port are
those of $\Gmat_{\Inp}$.

\begin{definition}\label{def:port-dilation}
Let $\Recv$ be a \tref{receiver instrument} with
\tref[port Gram matrices]{port Grams} $\{A_k\}_{k=1}^{K}$ on
$\mathbb{C}^{\Inp}$ (Definitions~\ref{def:instrument} and~\ref{def:gram}). A
\term[a finite ordered mode set with disjoint ports and an isometry from the
declared input modes into it whose Gram on each port is that port's Gram
matrix]{port dilation} of $\Recv$ is a triple
$(\mathbb{O};\mathbb{O}_1,\dots,\mathbb{O}_K;\Omega)$ of a finite mode set
$\mathbb{O}$ with a fixed order, pairwise disjoint subsets
$\mathbb{O}_1,\dots,\mathbb{O}_K$ of $\mathbb{O}$, its \emph{ports}, and a
linear map $\Omega:\mathbb{C}^{\Inp}\to\mathbb{C}^{\mathbb{O}}$ with
\begin{equation}\label{eq:port-dilation}
  \Omega^{\dagger}\Omega=\id_{\Inp}
  \qquad\text{and}\qquad
  \Omega^{\dagger}\Emd{\mathbb{O}_k}\Omega=A_k\quad(k=1,\dots,K),
\end{equation}
where $\Emd{\mathbb{O}_k}$ is the coordinate projector of
$\mathbb{C}^{\mathbb{O}}$ onto the modes of $\mathbb{O}_k$. The modes of
$\mathbb{O}$ in no port are its \emph{loss modes}, and the \emph{row} of a mode
$i\in\mathbb{O}$ is the row of $\Omega$ that $i$ indexes. The \emph{click
pattern} of an occupation vector $\vec m\in\mathbb{N}_0^{\mathbb{O}}$
(Definition~\ref{def:fock}) is
\begin{equation}\label{eq:click-pattern-of}
  S_{\vec m}\coloneqq\big\{k\in[K]:\ m_i\ge1\text{ for some }i\in\mathbb{O}_k\big\}.
\end{equation}
\end{definition}
